\documentclass[10pt]{article}
\usepackage[margin=1in]{geometry}
\usepackage{booktabs}
\usepackage{hyperref}
\usepackage{enumitem}
\usepackage{parskip}
\usepackage{microtype}
\usepackage{titlesec}

\titleformat{\section}{\normalfont\bfseries}{\thesection.}{0.5em}{}
\titlespacing{\section}{0pt}{1.2ex}{0.5ex}

\hypersetup{colorlinks=true, linkcolor=blue, urlcolor=blue, citecolor=blue}

\begin{document}

\begin{center}
  {\large\bfseries Rebuttal --- Reviewer 3}
\end{center}

We are grateful for the thoughtful and generous review. The suggestions are
all actionable and we address each in turn.

% ---------------------------------------------------------------
\section*{Central Figure}

We agree this would significantly help readers orient early. We will add a
summary figure (end of Section~1 or beginning of Section~3) covering:
(a)~the semantic MM~$\leftrightarrow$~Trader dialogue with a Brier-score
comparison; (b)~the LMSR mechanism with agent roles and price dynamics; and
(c)~a headline comparison of market vs.\ debate vs.\ voting across adversarial
conditions. This gives readers a visual anchor before the technical details.

% ---------------------------------------------------------------
\section*{Conditions Where the Mechanism Fails --- Moving to the Introduction}

This is good pedagogical advice. The limitations are currently confined to
Section~7, which means readers encounter them only after investing in the full
technical presentation. We will move a brief, honest summary of the key failure
modes into Section~1 or the paper's abstract:
(1)~the mechanism assumes adversarial agents are imperfect (i.e.\ not fully
rational wealth-maximisers);
(2)~performance degrades significantly in multi-choice settings due to liquidity
fragmentation;
(3)~very weak market makers may be unable to distinguish valid from deceptive
reasoning in the semantic variant.
Flagging these upfront sets appropriate expectations and, as the reviewer notes,
frees up working memory to engage with what the mechanism does well.

% ---------------------------------------------------------------
\section*{Related Literature}

We thank the reviewer for these pointers and will conduct a broader sweep before
the revision.

\textbf{Todasco (2025)~\cite{todasco2025}:} This work uses stake size as a
confidence signal from individual LLMs, showing that larger bets correlate with
higher accuracy ($\approx\!99\%$ at maximum stakes vs.\ $74\%$ at minimum
stakes). The conceptual connection to our work is real---both treat financial
commitment as a proxy for epistemic certainty---but the mechanisms are distinct.
Todasco extracts confidence from a \emph{single} agent through stake calibration;
our paper uses \emph{multi-agent} LMSR markets for collective truth elicitation
and adversarial robustness. The single-agent result is complementary: it suggests
that the financial penalty structure in our LMSR setting may have additional
confidence-calibrating effects beyond what we currently measure. We will position
this as related but distinct in Section~2.2.

\textbf{OpenForecaster~\cite{openforecaster}:} This project trains a
specialised 8B forecasting model optimised for Brier score and calibration using
GRPO reinforcement learning. It addresses the complementary question of
\emph{training} models to be well-calibrated, whereas our work addresses
\emph{coordinating} models that were not specifically trained for forecasting.
We will cite this in the related work as an orthogonal approach to the same
calibration goal.

% ---------------------------------------------------------------
\section*{Moderators of Mechanism Efficacy}

These are valuable questions we will address in an expanded discussion.
On \textbf{time horizon}: our simulations suggest wealth divergence is visible
within 20--30 episodes for typical adversarial effectiveness values; we will add
a figure showing this. On \textbf{external reward structure}: our results already
show ground-truth resolution outperforms consensus, consistent with the
reviewer's intuition; we will discuss this as a direction for future work. On
\textbf{capability differential}: Section~5's weak-to-strong experiments begin
to address this (GPT-4.1-nano successfully extracts signal from GPT-4.1-standard
under adversarial conditions); we will expand the discussion of where this breaks
down when the capability gap becomes extreme.

% ---------------------------------------------------------------
\section*{Title}

We appreciate the suggestion and agree the current title under-sells the
comparative result. We will consider a revised title that more directly signals
the head-to-head finding, while remaining consistent with the venue's conventions.

% ---------------------------------------------------------------
\section*{Revision Commitments}

\begin{itemize}[noitemsep, topsep=2pt, leftmargin=*]
  \item Add a central summary figure (mechanisms + headline comparison) in
        Section~1 or~3
  \item Move key failure modes to the introduction: adversarial imperfection
        assumption, MCQ degradation, semantic market maker capability threshold
  \item Add Todasco (2025) and OpenForecaster to Section~2.2 with appropriate
        positioning
  \item Conduct a broader related-work sweep for prediction-market-as-inference
        literature
  \item Expand Section~7 discussion of moderators: episode count, external
        reward structure, capability differential limits
  \item Consider a more descriptive title that reflects the comparative result
\end{itemize}

\end{document}
